\chapter{Reading the log}
\label{ch:log}

Experts answer ``what happened?'' from the log, not from memory. Every
session is one file of JSON lines, and everything you saw (and some things
you did not) is in it.

\section{Finding a session}

\begin{sh}
nodwyn sessions                          # id, last event, state, folder, first message
nodwyn replay SESSION                    # the transcript, rebuilt from the file
nodwyn --show-thinking replay SESSION    # with the model's thinking
\end{sh}

\begin{txt}
# session 20260930T201929Z-fd1fbe10 — demo ollama @ /home/you/garden (15 tools)

you> Summarize this folder
▸ fs-list {"path":"."}
▸ fs-read {"path":"README.md"}
◂ fs-list ok (37 bytes)
◂ fs-read ok (45 bytes)
## What's here
...
# end (2026-09-30T20:19:30Z)
\end{txt}

The file itself is
\file{\textasciitilde/.local/share/nodwyn/sessions/SESSION.jsonl}. Each line
has a \code{seq}, a timestamp, the session id, a \code{kind} and a
\code{payload}:

\begin{txt}
{"seq":6,"ts":"2026-09-30T20:19:29.770583Z","session":"20260930T201929Z-fd1fbe10",
 "kind":"guard.decision",
 "payload":{"args":"(\".\")","decision":"allow","origin":"model","prim":"fs-list"}}
\end{txt}

Beside it, \file{SESSION.mem} holds the session's memory (ADR 51): a card
for each whole turn and each compaction's note, and the words of the events
\code{index-search} ranks (ADR 50), kept and checked against the log the way
the index is, and made again whenever it does not match. And
\file{SESSION.idx} is the log's index (ADR 49): the session's
listing and where every 256th line starts, so that listing the session and
reading a page of its log do not read the whole log again as it grows. It
is derived from the log and checked against the log's line where it ends
before use; one that does not match is built again from the log, so it can
always be deleted. Resuming a session reads the index and the end of the
log: the session holds the events from the point its window starts (where a
request last cut the conversation at a lasting narrowing, a compaction's),
or from its last three requests or an unfinished plan's turn if those are
earlier, and its model, judge and limits come from the index.

\section{Human search and focused history}

The console's search (the sidebar box, or a task's menu) uses \code{conversation.search} over the
same word postings and turn cards as the CLI and model. Its anchor task's
recorded project bounds earlier scope, including archived and worktree tasks;
the model's earlier-tasks switch remains unchanged. Both new methods are
read-only and available on paired and read-only consoles. Ranked results
are grouped by task and turn, with event sequences, times and redacted
snippets. Failed source reads make the result explicitly incomplete.

A match opens a separate historical Transcript through \code{log.context}.
Indexed run and turn boundaries fix its range; the first page has at most
40 neighbouring records on either side. Earlier/later actions read up to
500 records inside that range using \code{log.read}. The match's folds
open and its rendered node is highlighted; historical choices and receipts
never mutate live interaction. Open task returns to the ordinary current
conversation. Query and scope survive reload and Back, and stale responses
cannot replace a newer view. The broader History panel, search settings,
semantic ranking, providers' embeddings and measurement corpus remain
separate work (ADR 51).

\section{The event kinds}

\begin{center}
\begin{tabularx}{\linewidth}{lX}
\toprule
\code{session.start}, \code{session.resume}, \code{session.end} & lifecycle; no \code{session.end} after a crash \\
\code{session.title}, \code{session.archive} & the task's name, the one you gave it or the one the model wrote from its first message (actor \code{model}), and whether it is archived (\adr{39}); the latest of each counts \\
\code{session.model} & the model the conversation talks to, chosen as it opened or moved to between turns (\adr{40}), with the window's recipe sized for it; the latest one counts \\
\code{worktree.create}, \code{worktree.restore}, \code{worktree.remove} & the task's own git worktree (\adr{34}): made, made again on resume, removed on close when nothing in it was unlanded \\
\code{user.msg}, \code{harness.msg} & what you sent; what a strategy added with \code{add-message} \\
\code{llm.req}, \code{llm.resp} & each model request (with its hash and recipe) and reply, with its usage: tokens, and cached tokens and cost when known \\
\code{tool.call}, \code{tool.result} & each call and its one result \\
\code{guard.decision} & every confirm and deny, and every allow of a tool or a host primitive, with checker, reason and answer \\
\code{egress.check}, \code{jev.req} & keyed requests: binding lookup; judge calls \\
\code{gate.plan}, \code{plan.step} & the plan gate's outcome; the model's marks on its steps \\
\code{context.narrow}, \code{context.fold}, \code{context.widen}, \code{result.fold} & window shaping \\
\code{context.compact} & a compaction: what set it off, the context's size before and after, the note \\
\code{turn.checkpoint}, \code{turn.end} & budget checkpoints; the durable end of a turn, with its outcome \\
\code{run.state}, \code{run.turn}, \code{run.checkpoint} & runs (\code{nodwyn run}) \\
\code{eval.req}, \code{eval.result} & your \code{:eval}s \\
\code{redact.event} & a redaction: where and which class, never the value \\
\code{cred.register}, \code{cred.refresh}, \code{binding.change}, \code{approval.change}, \code{setting.change} & keys, policy, and a setting you chose at a prompt \\
\code{file.version}, \code{attach}, \code{doc.upload} & kept versions; attachments; documents sent out \\
\code{summary}, \code{summary.drop} & reviewed summaries \\
\code{note} & \code{log-note} from a strategy, and the harness's own notes: a retried model request, a judged approval, a resume \\
\bottomrule
\end{tabularx}
\end{center}

\section{\texttt{jq} recipes}

\begin{sh}
F=~/.local/share/nodwyn/sessions/SESSION.jsonl

jq -c '[.seq,.kind]' $F                                   # the outline
jq -c 'select(.kind=="guard.decision") | .payload
       | {prim,origin,decision,checker,reason}' $F        # every decision
jq -c 'select(.kind=="guard.decision" and .payload.decision!="allow")
       | .payload' $F                                     # only what was stopped
jq -r 'select(.kind=="tool.result" and .payload.ok==false)
       | "\(.seq) \(.payload.name): \(.payload.error)"' $F # failures
jq -r 'select(.kind=="note") | .payload.text' $F           # notes
jq -c 'select(.kind=="llm.resp") | .payload.usage
       | {prompt_tokens,cached_tokens,cost}' $F           # what each request cost
\end{sh}

\begin{tip}
A \code{guard.decision} with \code{origin} \code{harness} came from the
loop's own strategy (the planner, compaction, a sub-agent's plumbing);
\code{model} is the model's call or something a tool it called did;
\code{user} is you. When a decision surprises you, \code{checker} tells you
which rule made it and \code{reason} tells you why.
\end{tip}

\section{Exactly what the model saw}

\code{llm.req} records a hash and the recipe of the window, not the whole
prompt. Rebuild the request, byte for byte, from the log before it:

\begin{sh}
nodwyn replay SESSION --window 3
\end{sh}

\begin{txt}
{
  "Model": "demo",
  "Messages": [
    {
      "role": "system",
      "content": "You are nodwyn, a coding agent working inside the repository at ...
\end{txt}

The command checks the rebuilt request against the recorded hash and
refuses to print a mismatch. It never calls the model or runs a tool. This
is the tool for ``why did the model do \emph{that}?'': read what it was
actually shown.

\section{Summaries you write}

For a long conversation you want to keep working in, you can replace a
finished stretch with a summary you wrote and reviewed:

\begin{sh}
nodwyn summary --from 12 --through 41 --text summary.txt SESSION   # add
nodwyn summary SESSION                                              # list
nodwyn summary --drop SUMMARY_SEQ SESSION                           # revoke
\end{sh}

The span must start at a \code{user.msg} and end at a finished reply. The
window then shows the summary, labelled as a revocable summary, in place of
that span; the original events stay in the log, and dropping the summary
brings them back. The latest two user turns are never replaced. Close the
session before editing its summaries.

\section{When something is wrong}

\begin{center}
\begin{tabularx}{\linewidth}{lX}
\toprule
\textbf{symptom} & \textbf{look at} \\
\midrule
\code{no control model is configured} & add the \code{[control]} table it prints \\
\code{control model \ldots\ unreachable} & the \code{url}; for Ollama, \code{curl URL/api/version} \\
first turn takes 10--60 s & a cold prompt cache, or another client holds the model server \\
every request is slow & the prefix cache misses: look for trims (the \code{elided} count jumping) or folds in \code{:log} \\
a hosted model costs more than expected & the usage line's cached share (a low one is the same cache miss) and \code{:budget} in the terminal client; prices in the model's table if no cost shows \\
the model loops on a denied call & the progress guard will stop it; read the \code{guard.decision}s first \\
\code{hard stop: tool\_calls budget reached} & the turn ended; ask for something smaller, or raise \code{-{}-max-tool-calls} \\
a tool call errored & \code{:log 20}: the \code{tool.result} has the exact error the model saw \\
a file changed without a confirm & by default writes in the folder go ahead: the decision's reason says \code{approval.writes}, or names the standing approval that covered it; git or \code{nodwyn undo} puts it back \\
a command ran without asking you & the decision's \code{actor}: \code{judge} (the judge allowed it; a \code{note} has its reason) or \code{setting} (\code{[approval] commands = "allow"}); otherwise \code{:allow} lists the standing approval \\
a strategy seems ignored & \code{:eval agent-turn} prints which procedure is installed; \code{(reload)} returns notes \\
\bottomrule
\end{tabularx}
\end{center}
